Năm 1985, David Deutsch đã đặt ra một câu hỏi tưởng chừng triết học: liệu có tồn tại một máy tính vật lý có thể mô phỏng hiệu quả mọi quá trình vật lý trong tự nhiên không? Câu trả lời mà ông đưa ra rằng, máy tính đó phải hoạt động theo các nguyên lý của cơ học lượng tử. Từ đó đã đặt nền móng cho một lĩnh vực hoàn toàn mới: tính toán lượng tử (quantum computation). Bảy năm sau, Deutsch và Jozsa đã phát triển một thuật toán để chứng minh một máy tính lượng tử có thể giải một bài toán cụ thể chỉ trong một bước duy nhất, trong khi bất kỳ thuật toán cổ điển tất định nào cũng cần đến hàng triệu bước khi kích thước đầu vào tăng lên.

Bốn thập kỷ sau, tính toán lượng tử đang chuyển dịch mạnh mẽ từ lý thuyết sang hiện thực. IBM, Google và IonQ lần lượt công bố các hệ thống vượt mốc 1000 qubit trong giai đoạn 2023–2025, đánh dấu bước chuyển từ thí nghiệm phòng lab sang nền tảng có thể tiếp cận qua đám mây. Tại Việt Nam, sự quan tâm đến lĩnh vực này đang gia tăng cùng với các chương trình đầu tư vào nhân lực nghiên cứu lượng tử, nhưng tài liệu tiếng Việt kết nối chặt chẽ giữa nền tảng toán học và lập trình lượng tử thực tiễn vẫn còn rất hạn chế.

Để hiểu tại sao tính toán lượng tử có thể vượt trội so với tính toán cổ điển, cần phải hiểu nền tảng toán học đằng sau nó. Máy tính cổ điển xử lý thông tin dưới dạng bit nhận giá trị 0 hoặc 1 -- giống như công tắc chỉ có thể tắt hoặc bật. Máy tính lượng tử thay thế bit bằng qubit (quantum bit), có thể tồn tại đồng thời ở cả hai trạng thái cho đến khi được đo, gọi là chồng chất lượng tử (quantum superposition). Xác suất của kết quả đo không phải là một con số tùy ý mà được xác định bởi biên độ xác suất (probability amplitude) trong không gian Hilbert phức, các biên độ này có thể triệt tiêu nhau như sóng, tạo ra hiện tượng giao thoa lượng tử (quantum interference) mà xác suất cổ điển không thể giải thích được. Chính yếu tố đó là cơ sở cho thuật toán Deutsch-Jozsa. Nền tảng toán học đầy đủ của lý thuyết này bao gồm: không gian Hilbert, toán tử tuyến tính, hệ tiên đề và quy tắc Born đều được trình bày có tính hệ thống trong các tài liệu về xác suất lượng tử và tính toán lượng tử. Tuy nhiên, các hướng dẫn lập trình lượng tử thực tiễn lại hiếm khi giải thích cơ sở đó. Điều đó tạo ra một rào cản lớn, khiến người học chỉ có thể thao tác lập trình mà thiếu đi sự hiểu biết sâu sắc về nguyên lý vận hành cốt lõi của các mạch lượng tử.

Mục tiêu của khóa luận này là xây dựng một trình bày toán học về cơ sở xác suất lượng tử từ không gian Hilbert và tiên đề cơ học lượng tử đến lý thuyết phép đo,  từ đó phân tích thuật toán Deutsch-Jozsa và kiểm chứng qua mô phỏng số với Python và Qiskit. Thuật toán Deutsch-Jozsa được chọn vì nó là đối tượng lý tưởng để minh họa trọn vẹn các khái niệm cốt lõi bao gồm: chồng chất, giao thoa, oracle lượng tử (quantum oracle) và phép đo lượng tử. Đây là mô hình phù hợp để phân tích chi tiết và trình bày một cách hệ thống trong khuôn khổ của một khóa luận tốt nghiệp đại học.

Khóa luận được tổ chức thành ba chương chính.
\begin{itemize}
    \item Chương 1 xây dựng nền tảng toán học của cơ học lượng tử, bao gồm không gian Hilbert, các lớp toán tử tuyến tính quan trọng, hệ tiên đề lượng tử.
    \item Chương 2 xây dựng mô hình tính toán lượng tử: qubit, các cổng lượng tử phổ biến và cấu trúc mạch.
    \item Chương 3 phân tích thuật toán Deutsch-Jozsa và kiểm chứng qua mô phỏng số.
\end{itemize}  

Trong quá trình thực hiện khóa luận, có một số nội dung được sử dụng các công cụ trí tuệ nhân tạo như một trợ lý soạn thảo và tra cứu: hỗ trợ diễn đạt lại các nội dung được dịch ra từ tài liệu tham khảo, gợi ý cấu trúc trình bày và kiểm tra tính nhất quán của ký hiệu toán học. Bản thân em cam kết toàn bộ nội dung toán học, diễn giải lập luận và kết quả mô phỏng đều được tự xây dựng, tổng hợp và chịu trách nhiệm.